%!TEX root = ../memoria.tex
% =============================================================================
%  CAPÍTULO 2 - ESTADO DEL ARTE (recortado, sin FaceTime/Meet, sin guiones)
% =============================================================================

Este capítulo sitúa \emph{ConectaMayor} en el contexto actual y justifica
la idea del proyecto. Para ello se analizan cuatro aplicaciones
representativas, con enfoques distintos para abordar la comunicación
familiar y la inclusión digital de las personas mayores. El análisis se
agrupa en tres bloques: la mensajería generalista, las soluciones
específicas para personas mayores y las plataformas de compartición
visual. Termina con una comparativa que recoge las conclusiones
principales.


\section{Aplicaciones de mensajería generalista\label{SEC:EA_GENERALISTAS}}

Las aplicaciones de mensajería generalista son, con diferencia, las más
utilizadas por las personas mayores que mantienen contacto digital con su
familia, aunque suelen recurrir a ellas por necesidad más que por
comodidad.

\subsection{WhatsApp}

WhatsApp \cite{whatsapp} es la aplicación de mensajería más extendida en
España. Permite enviar mensajes, audios, fotografías y vídeos, hacer
llamadas y videollamadas y crear grupos. Esa abundancia de funciones, que
para el público general es una ventaja, se convierte en un problema para
el usuario mayor. La interfaz reúne en una misma pantalla mensajes
individuales, grupos, llamadas, estados, canales y comunidades, sin que el
usuario pueda decidir con facilidad qué quiere ver y qué no. WhatsApp
favorece la comunicación, pero dentro de un entorno pensado para la
población general y poco adaptado al usuario mayor.


\section{Soluciones específicas para personas mayores\label{SEC:EA_ESPECIFICAS}}

Frente a las plataformas generalistas existe un conjunto reducido de
soluciones diseñadas específicamente para personas mayores. Son las
referencias más directas para \emph{ConectaMayor}, ya que comparten su
punto de partida: reconocer que las herramientas convencionales no cubren
bien las necesidades de este colectivo y proponer un entorno pensado para
ellas desde el inicio.

\subsection{GrandPad}

GrandPad \cite{grandpad} es probablemente la solución más conocida a nivel
internacional. Es una tableta de hardware específico, con un sistema
operativo y aplicaciones diseñadas íntegramente para personas mayores:
iconos grandes, un único nivel de navegación y un círculo familiar cerrado
en el que solo los allegados autorizados pueden contactar con el usuario.
Incluye llamadas, videollamadas, mensajes, fotografías, juegos y un
asistente humano accesible con un botón.

El planteamiento es acertado, pero tiene dos limitaciones que lo alejan de
muchas familias. La primera es económica, ya que se comercializa por
suscripción mensual, un coste difícil de asumir para una parte importante
de los hogares. La segunda es de modelo, porque al ser un dispositivo
dedicado obliga a adquirir un equipo nuevo en lugar de aprovechar uno que
la persona ya tenga.

\subsection{Mayores Conectados}

Mayores Conectados \cite{mayoresconectados} es un proyecto argentino que
parte del mismo diagnóstico que este trabajo: fomentar la relación
intergeneracional mediante la tecnología, a través de actividades
reconocibles para la persona mayor. La plataforma reúne juegos y
actividades para que abuelos y nietos las hagan juntos, con interfaces
simplificadas.

La diferencia principal con \emph{ConectaMayor} es el enfoque. Mayores
Conectados se orienta al ocio compartido, mientras que este proyecto se
centra en la comunicación cotidiana, por lo que son enfoques
complementarios. Una limitación del proyecto argentino es que parte de sus
juegos están alojados en plataformas externas, lo que rompe la sensación
de entorno único y puede generar desconfianza al saltar entre dominios.


\section{Plataformas de compartición visual\label{SEC:EA_VISUALES}}

El tercer bloque agrupa las aplicaciones cuyo valor principal es compartir
contenido visual, sobre todo fotografías. \emph{ConectaMayor} incorpora un
módulo de galería con esa misma orientación, por lo que conviene analizar
el referente más evidente.

\subsection{Instagram}

Instagram \cite{instagram} es una de las plataformas de referencia para
compartir fotografías, y fue una inspiración directa para el módulo de
galería. La idea inicial del proyecto era ofrecer una especie de Instagram
familiar y privado. Su cuadrícula de imágenes con comentarios ha
establecido un patrón visual que casi todo el mundo reconoce, incluidas
muchas personas mayores que la han visto a través de algún familiar.

Sin embargo, Instagram no encaja en el uso familiar que plantea este
proyecto por dos motivos. El primero es su naturaleza abierta, ya que,
aunque permite cuentas privadas, está pensada para la difusión pública y
la interacción con desconocidos, una superficie de exposición que no
encaja con el caso de una familia. El segundo es la saturación de su
interfaz, que mezcla la cuadrícula personal con contenido recomendado,
mensajería, vídeos cortos y compras. Para una persona mayor, eso convierte
algo sencillo como ver las fotos que comparte su familia en una
experiencia confusa. \emph{ConectaMayor} retoma la idea visual de
Instagram y la lleva a un entorno privado, sin contenido externo y
limitado al grupo familiar.


\section{Análisis comparativo\label{SEC:EA_COMPARATIVA}}

La Tabla~\ref{TAB:EA_COMPARATIVA} resume las características más relevantes
de las aplicaciones analizadas frente a \emph{ConectaMayor}. Los criterios
elegidos son el diseño específico para mayores, el ámbito privado, la
gratuidad, la interfaz unificada y la cobertura funcional, y recogen los
aspectos sobre los que se articula el resto de la memoria.

\begin{table}[Comparativa de aplicaciones]{TAB:EA_COMPARATIVA}{Comparativa de las aplicaciones analizadas frente a \emph{ConectaMayor}.}
  \small
  \begin{tabular}{lccccc}
    \toprule
    \textbf{Aplicación} & \textbf{Diseño} & \textbf{Ámbito} & \textbf{Gratuita} & \textbf{Interfaz} & \textbf{Cobertura} \\
                        & \textbf{para mayores} & \textbf{privado} & & \textbf{unificada} & \textbf{funcional} \\
    \midrule
    WhatsApp             & No  & Parcial & Sí  & No  & Media \\
    GrandPad             & Sí  & Sí      & No  & Sí  & Alta \\
    Mayores Conectados   & Sí  & Sí      & Sí  & No  & Media \\
    Instagram            & No  & No      & Sí  & No  & Media \\
    \midrule
    \textbf{ConectaMayor} & \textbf{Sí} & \textbf{Sí} & \textbf{Sí} & \textbf{Sí} & \textbf{Alta} \\
    \bottomrule
  \end{tabular}
\end{table}

Del análisis se desprenden tres conclusiones que justifican el proyecto.
La primera es que ninguna de las soluciones gratuitas analizadas combina a
la vez un diseño pensado para personas mayores, un ámbito estrictamente
privado y una cobertura funcional amplia. WhatsApp es gratuita pero no
está adaptada. GrandPad está adaptado y es privado, pero tiene un coste
recurrente. Mayores Conectados está adaptado y es gratuito, pero se apoya
en plataformas externas. Instagram es gratuito, pero ni está adaptado ni
es privado.

La segunda es que las soluciones existentes abordan dimensiones aisladas
del problema, como la mensajería, la compartición visual o el ocio, pero
no integran un conjunto coherente de servicios en torno a las necesidades
cotidianas de la persona mayor. Esa fragmentación obliga a manejar varias
aplicaciones para tareas distintas, lo que aumenta la carga cognitiva y la
probabilidad de abandono.

La tercera es que ninguna implementa un sistema de roles familiares como
el de \emph{ConectaMayor}, en el que un familiar puede asumir la gestión,
creando recordatorios, manteniendo los contactos o supervisando el uso,
sin invadir el espacio personal de la persona mayor. Este es uno de los
rasgos diferenciales del proyecto y se desarrolla en el
Capítulo~\ref{CAP:DISENO}.

\emph{ConectaMayor} no pretende reemplazar a ninguna de estas
herramientas, sino ocupar un espacio que hoy no cubre ninguna: una
aplicación web gratuita, adaptada a las personas mayores, privada por
diseño, con interfaz unificada y con un sistema de roles que reconoce la
naturaleza colaborativa del uso familiar de la tecnología.
